\begin{bt}%[9D6H18-2]
% nguon: Cụm 10 – Đề tham khảo Tuyển sinh 10 năm học 2026-2027 TP.HCM – Đề số 30
Cho hàm số $y = -\dfrac{1}{2}x^2$ có đồ thị là parabol $(P)$.
\begin{enumerate}[a)]
    \item Vẽ đồ thị $(P)$ trên mặt phẳng tọa độ $Oxy$.
    \item Tìm tọa độ các điểm $M$ (khác gốc tọa độ) thuộc $(P)$ sao cho điểm $M$ cách đều hai trục tọa độ.
\end{enumerate}
\loigiai{
\begin{enumerate}[a)]
    \item Bảng giá trị:
    \begin{center}
    \begin{tabular}{|c|c|c|c|c|c|}
    \hline
    $x$ & $-4$ & $-2$ & $0$ & $2$ & $4$ \\ \hline
    $y = -\dfrac{1}{2}x^2$ & $-8$ & $-2$ & $0$ & $-2$ & $-8$ \\ \hline
    \end{tabular}
    \end{center}
    Đồ thị parabol $(P)$ đi qua gốc tọa độ $O(0;0)$ và các điểm $(-4;-8)$, $(-2;-2)$, $(2;-2)$, $(4;-8)$.
    \begin{center}
    \begin{tikzpicture}[>=stealth,scale=0.7]
        \draw[->] (-5,0) -- (5,0) node[below] {$x$};
        \draw[->] (0,-9) -- (0,1.5) node[left] {$y$};
        \node[above left] at (0,0) {$O$};
        \draw[dashed] (-4,0) node[above] {$-4$} -- (-4,-8) -- (4,-8) -- (4,0) node[above] {$4$};
        \draw[dashed] (-2,0) node[above] {$-2$} -- (-2,-2) -- (2,-2) -- (2,0) node[above] {$2$};
        \draw[dashed] (-4,-8) -- (0,-8) node[right] {$-8$};
        \draw[dashed] (-2,-2) -- (0,-2) node[right] {$-2$};
        \draw[thick,smooth,samples=100,domain=-4.2:4.2] plot(\x,{-0.5*(\x)^2});
        \fill (-4,-8) circle (1.5pt) (-2,-2) circle (1.5pt) (0,0) circle (1.5pt) (2,-2) circle (1.5pt) (4,-8) circle (1.5pt);
    \end{tikzpicture}
    \end{center}
    \item Vì điểm $M(x;y)$ cách đều hai trục tọa độ nên $|y| = |x| \iff y = x$ hoặc $y = -x$.
    \begin{itemize}
        \item Trường hợp 1: $y = x$.\\
        Ta có phương trình hoành độ giao điểm:
        $$x = -\dfrac{1}{2}x^2 \iff \dfrac{1}{2}x^2 + x = 0 \iff x(x + 2) = 0 \iff \left[\begin{aligned} &x = 0\text{ (loại vì } M \ne O\text{)} \\ &x = -2\text{ (nhận)}. \end{aligned}\right.$$
        Với $x = -2 \implies y = -2$.
        \item Trường hợp 2: $y = -x$.\\
        Ta có phương trình hoành độ giao điểm:
        $$-x = -\dfrac{1}{2}x^2 \iff \dfrac{1}{2}x^2 - x = 0 \iff x(x - 2) = 0 \iff \left[\begin{aligned} &x = 0\text{ (loại vì } M \ne O\text{)} \\ &x = 2\text{ (nhận)}. \end{aligned}\right.$$
        Với $x = 2 \implies y = -2$.
    \end{itemize}
    Vậy tọa độ các điểm cần tìm là $(-2;-2)$ và $(2;-2)$.
\end{enumerate}
}
\end{bt}

\begin{bt}%[9D6H20-2]
% nguon: Cụm 10 – Đề tham khảo Tuyển sinh 10 năm học 2026-2027 TP.HCM – Đề số 30
Cho phương trình $x^2 - 5x - 3 = 0$ có hai nghiệm $x_1, x_2$. Không giải phương trình, hãy tính giá trị biểu thức:
$$A = (x_2 + 5)(x_2^2 - 3) + x_1^2 + x_2^2 + 50x_1.$$
\loigiai{
Phương trình $x^2 - 5x - 3 = 0$ có:
$$\Delta = (-5)^2 - 4 \cdot 1 \cdot (-3) = 25 + 12 = 37 > 0.$$
Vì $\Delta > 0$ nên phương trình có hai nghiệm phân biệt $x_1, x_2$.\\
Theo định lý Vi-ét, ta có:
$$\begin{cases} x_1 + x_2 = 5 \\ x_1 x_2 = -3. \end{cases}$$
Ta có $x_1^2 + x_2^2 = (x_1 + x_2)^2 - 2x_1 x_2 = 5^2 - 2 \cdot (-3) = 25 + 6 = 31$.\\
Vì $x_2$ là nghiệm của phương trình nên $x_2^2 - 5x_2 - 3 = 0 \implies x_2^2 - 3 = 5x_2$.\\
Thay vào biểu thức $A$, ta được:
$$\begin{aligned}
A &= (x_2 + 5) \cdot 5x_2 + x_1^2 + x_2^2 + 50x_1 \\
&= 5x_2^2 + 25x_2 + x_1^2 + x_2^2 + 50x_1 \\
&= 5(5x_2 + 3) + 25x_2 + x_1^2 + x_2^2 + 50x_1 \\
&= 25x_2 + 15 + 25x_2 + x_1^2 + x_2^2 + 50x_1 \\
&= 50x_1 + 50x_2 + (x_1^2 + x_2^2) + 15 \\
&= 50(x_1 + x_2) + (x_1^2 + x_2^2) + 15.
\end{aligned}$$
Thay $x_1 + x_2 = 5$ và $x_1^2 + x_2^2 = 31$ vào, ta được:
$$A = 50 \cdot 5 + 31 + 15 = 250 + 46 = 296.$$
}
\end{bt}

\begin{bt}%[9D8H26-3]
% nguon: Cụm 10 – Đề tham khảo Tuyển sinh 10 năm học 2026-2027 TP.HCM – Đề số 30
Số học sinh nam và nữ của $6$ lớp khối $9$ tại một trường THCS được thống kê như sau:
\begin{center}
\begin{tabular}{|l|c|c|c|c|c|c|}
\hline
\textbf{Lớp} & \textbf{9A} & \textbf{9B} & \textbf{9C} & \textbf{9D} & \textbf{9E} & \textbf{9F} \\ \hline
Số học sinh Nam & $21$ & $13$ & $27$ & $20$ & $25$ & $14$ \\ \hline
Số học sinh Nữ & $19$ & $29$ & $15$ & $18$ & $20$ & $30$ \\ \hline
\end{tabular}
\end{center}
\begin{enumerate}[a)]
    \item Lớp có tổng số học sinh nhiều nhất là lớp nào?
    \item Chọn ngẫu nhiên một lớp trong $6$ lớp đó. Tính xác suất của các biến cố sau:
    \begin{itemize}
        \item $A$: ``Lớp được chọn có số học sinh nam nhiều hơn số học sinh nữ''.
        \item $B$: ``Lớp được chọn có tổng số học sinh không quá $44$ học sinh''.
    \end{itemize}
\end{enumerate}
\loigiai{
\begin{enumerate}[a)]
    \item Tổng số học sinh của các lớp lần lượt là:
    \begin{itemize}
        \item Lớp 9A: $21 + 19 = 40$ (học sinh).
        \item Lớp 9B: $13 + 29 = 42$ (học sinh).
        \item Lớp 9C: $27 + 15 = 42$ (học sinh).
        \item Lớp 9D: $20 + 18 = 38$ (học sinh).
        \item Lớp 9E: $25 + 20 = 45$ (học sinh).
        \item Lớp 9F: $14 + 30 = 44$ (học sinh).
    \end{itemize}
    Vậy lớp có tổng số học sinh nhiều nhất là lớp 9E ($45$ học sinh).
    \item Số phần tử của không gian mẫu là $n(\Omega) = 6$.
    \begin{itemize}
        \item Các lớp có số học sinh nam nhiều hơn số học sinh nữ là: 9A, 9C, 9D, 9E.\\
        Số kết quả thuận lợi cho biến cố $A$ là $n(A) = 4$.\\
        Xác suất của biến cố $A$ là:
        $$P(A) = \dfrac{n(A)}{n(\Omega)} = \dfrac{4}{6} = \dfrac{2}{3}.$$
        \item Các lớp có tổng số học sinh không quá $44$ học sinh ($\le 44$) là: 9A, 9B, 9C, 9D, 9F.\\
        Số kết quả thuận lợi cho biến cố $B$ là $n(B) = 5$.\\
        Xác suất của biến cố $B$ là:
        $$P(B) = \dfrac{n(B)}{n(\Omega)} = \dfrac{5}{6}.$$
    \end{itemize}
\end{enumerate}
}
\end{bt}

\begin{bt}%[9D6V21-3]
% nguon: Cụm 10 – Đề tham khảo Tuyển sinh 10 năm học 2026-2027 TP.HCM – Đề số 30
Cô Dung có một mảnh đất hình chữ nhật với chiều dài $28\text{ m}$ và chiều rộng $21\text{ m}$. Cô dự định xây một ngôi nhà hình chữ nhật ở một góc của mảnh đất, phần còn lại dùng làm sân vườn bao quanh. Sân vườn phía trên rộng $x\text{ (m)}$ và sân vườn phía bên phải rộng $(x + 3)\text{ (m)}$ với $0 < x < 21$.
\begin{enumerate}[a)]
    \item Viết biểu thức theo $x$ biểu thị diện tích mảnh đất dành để xây nhà.
    \item Tìm $x$, biết diện tích mảnh đất dành để xây nhà là $192\text{ m}^2$.
\end{enumerate}
\loigiai{
\begin{enumerate}[a)]
    \item Chiều rộng của phần đất xây nhà là:
    $$21 - x\text{ (m)}.$$
    Chiều dài của phần đất xây nhà là:
    $$28 - (x + 3) = 25 - x\text{ (m)}.$$
    Biểu thức biểu diễn diện tích mảnh đất xây nhà theo $x$ là:
    $$S = (25 - x)(21 - x) = x^2 - 46x + 525\text{ (m}^2\text{)}.$$
    \item Theo đề bài, diện tích phần đất xây nhà là $192\text{ m}^2$, nên ta có phương trình:
    $$x^2 - 46x + 525 = 192 \iff x^2 - 46x + 333 = 0.$$
    Ta có $\Delta' = (-23)^2 - 1 \cdot 333 = 529 - 333 = 196 > 0$, $\sqrt{\Delta'} = 14$.\\
    Phương trình có hai nghiệm:
    $$x_1 = 23 - 14 = 9\text{ (thỏa mãn } 0 < x < 21\text{)},$$
    $$x_2 = 23 + 14 = 37\text{ (loại vì } x > 21\text{)}.$$
    Vậy $x = 9\text{ m}$.
\end{enumerate}
}
\end{bt}

\begin{bt}%[9H0V31-3]
% nguon: Cụm 10 – Đề tham khảo Tuyển sinh 10 năm học 2026-2027 TP.HCM – Đề số 30
Một đống cát có dạng hình nón với đường kính đáy là $28\text{ m}$ và chiều cao là $50\text{ dm}$.
\begin{enumerate}[a)]
    \item Tính thể tích của đống cát trên với độ chính xác $d = 0{,}005$ (làm tròn kết quả đến chữ số thập phân thứ hai). Biết công thức tính thể tích hình nón là $V = \dfrac{1}{3}\pi R^2 h$.
    \item Để vận chuyển $\dfrac{1}{5}$ đống cát đó đến công trình xây dựng (sử dụng kết quả thể tích đã làm tròn ở câu a), người ta dùng xe tải có thùng dạng hình hộp chữ nhật kích thước $3{,}8\text{ m} \times 2{,}1\text{ m} \times 0{,}8\text{ m}$. Trong mỗi chuyến xe, thùng xe có thể chở lượng cát nhiều hơn thể tích thực của nó là $8\%$. Hỏi cần ít nhất bao nhiêu chuyến xe tải để chở hết lượng cát đó?
\end{enumerate}
\loigiai{
\begin{enumerate}[a)]
    \item Bán kính đáy của đống cát hình nón là:
    $$R = 28 : 2 = 14\text{ (m)}.$$
    Đổi: $50\text{ dm} = 5\text{ m}$.\\
    Thể tích của đống cát là:
    $$V = \dfrac{1}{3}\pi R^2 h = \dfrac{1}{3}\pi \cdot 14^2 \cdot 5 = \dfrac{980\pi}{3} \approx 1026{,}2536\text{ (m}^3\text{)}.$$
    Làm tròn với độ chính xác $d = 0{,}005$ (đến hàng phần trăm), ta được:
    $$V \approx 1026{,}25\text{ (m}^3\text{)}.$$
    \item Lượng cát cần vận chuyển là:
    $$V_{\text{cần}} = \dfrac{1}{5} \cdot 1026{,}25 = 205{,}25\text{ (m}^3\text{)}.$$
    Thể tích thực của thùng xe tải là:
    $$V_{\text{thực}} = 3{,}8 \cdot 2{,}1 \cdot 0{,}8 = 6{,}384\text{ (m}^3\text{)}.$$
    Lượng cát tối đa mỗi chuyến xe tải chở được là:
    $$V_{\text{chở}} = 6{,}384 + 6{,}384 \cdot 8\% = 6{,}384 \cdot 1{,}08 = 6{,}89472\text{ (m}^3\text{)}.$$
    Số chuyến xe cần thiết để vận chuyển hết số cát đó là:
    $$\dfrac{V_{\text{cần}}}{V_{\text{chở}}} = \dfrac{205{,}25}{6{,}89472} \approx 29{,}769\text{ (chuyến)}.$$
    Vì số chuyến xe phải là số tự nhiên nên cần ít nhất $30$ chuyến xe tải.
\end{enumerate}
}
\end{bt}

\begin{bt}%[9D1H3-4]
% nguon: Cụm 10 – Đề tham khảo Tuyển sinh 10 năm học 2026-2027 TP.HCM – Đề số 30
Một người đi xe đạp từ $A$ đến $B$ với vận tốc $12\text{ km/h}$ rồi đi tiếp từ $B$ đến $C$ với vận tốc $6\text{ km/h}$, tổng thời gian đi từ $A$ đến $C$ hết $75$ phút. Khi về, người đó đi từ $C$ đến $B$ với vận tốc $8\text{ km/h}$ rồi đi từ $B$ đến $A$ với vận tốc $4\text{ km/h}$, tổng thời gian về từ $C$ đến $A$ hết $1$ giờ $30$ phút. Tính chiều dài các quãng đường $AB$ và $BC$.
\loigiai{
Gọi $x\text{ (km)}$ là chiều dài quãng đường $AB$ và $y\text{ (km)}$ là chiều dài quãng đường $BC$ ($x > 0, y > 0$).\\
Đổi: $75\text{ phút} = \dfrac{5}{4}\text{ giờ}$; $1\text{ giờ } 30\text{ phút} = \dfrac{3}{2}\text{ giờ}$.
\begin{itemize}
    \item Khi đi từ $A$ đến $C$:
    Thời gian đi đoạn $AB$ là $\dfrac{x}{12}\text{ (giờ)}$, thời gian đi đoạn $BC$ là $\dfrac{y}{6}\text{ (giờ)}$.\\
    Theo bài ra ta có phương trình:
    $$\dfrac{x}{12} + \dfrac{y}{6} = \dfrac{5}{4} \quad (1)$$
    \item Khi về từ $C$ đến $A$:
    Thời gian đi đoạn $CB$ là $\dfrac{y}{8}\text{ (giờ)}$, thời gian đi đoạn $BA$ là $\dfrac{x}{4}\text{ (giờ)}$.\\
    Theo bài ra ta có phương trình:
    $$\dfrac{x}{4} + \dfrac{y}{8} = \dfrac{3}{2} \quad (2)$$
\end{itemize}
Từ $(1)$ và $(2)$ ta có hệ phương trình:
$$\begin{cases} \dfrac{x}{12} + \dfrac{y}{6} = \dfrac{5}{4} \\ \dfrac{x}{4} + \dfrac{y}{8} = \dfrac{3}{2} \end{cases} \iff \begin{cases} x + 2y = 15 \\ 2x + y = 12 \end{cases} \iff \begin{cases} x = 3 \\ y = 6 \end{cases}\text{ (thỏa mãn điều kiện)}.$$
Vậy chiều dài quãng đường $AB$ là $3\text{ km}$ và quãng đường $BC$ là $6\text{ km}$.
}
\end{bt}

\begin{bt}%[9H9C29-3]
% nguon: Cụm 10 – Đề tham khảo Tuyển sinh 10 năm học 2026-2027 TP.HCM – Đề số 30
Cho tam giác $ABC$ nhọn ($AB < AC$) nội tiếp đường tròn tâm $O$. Các đường cao $AD, BE, CF$ của tam giác $ABC$ cắt nhau tại trực tâm $H$. Gọi $M, N$ lần lượt là trung điểm của $BC$ và $AH$.
\begin{enumerate}[a)]
    \item Chứng minh tứ giác $BCEF$ và tứ giác $AEHF$ nội tiếp. Từ đó suy ra $MN \perp EF$.
    \item Gọi $K$ là giao điểm của $EF$ và $BC$. Chứng minh $KB \cdot KC = KF \cdot KE$.
    \item Biết $\widehat{BAC} = 60^\circ$ và $BC = 16\text{ cm}$. Tính bán kính đường tròn ngoại tiếp tam giác $AEF$.
\end{enumerate}
\loigiai{
\begin{center}
\begin{tikzpicture}[scale=0.95]
    \coordinate (A) at (1,4.5);
    \coordinate (B) at (-3,0);
    \coordinate (C) at (4,0);
    \coordinate (D) at ($(B)!(A)!(C)$);
    \coordinate (E) at ($(A)!(B)!(C)$);
    \coordinate (F) at ($(A)!(C)!(B)$);
    \coordinate (H) at (intersection of B--E and C--F);
    \coordinate (M) at ($(B)!0.5!(C)$);
    \coordinate (N) at ($(A)!0.5!(H)$);
    \coordinate (K) at (intersection of E--F and B--C);
    \draw[thick] (A) -- (B) -- (C) -- cycle;
    \draw[thick] (B) -- (E);
    \draw[thick] (C) -- (F);
    \draw[dashed] (A) -- (D);
    \draw[thick] (K) -- (E);
    \draw[thick] (K) -- (B);
    \draw[thick,blue] (M) -- (N);
    \node[above] at (A) {$A$};
    \node[below left] at (B) {$B$};
    \node[below right] at (C) {$C$};
    \node[below] at (D) {$D$};
    \node[above right] at (E) {$E$};
    \node[left] at (F) {$F$};
    \node[above right] at (H) {$H$};
    \node[below] at (M) {$M$};
    \node[above right] at (N) {$N$};
    \node[below left] at (K) {$K$};
    \fill (A) circle (1.5pt) (B) circle (1.5pt) (C) circle (1.5pt) (D) circle (1.5pt) (E) circle (1.5pt) (F) circle (1.5pt) (H) circle (1.5pt) (M) circle (1.5pt) (N) circle (1.5pt) (K) circle (1.5pt);
\end{tikzpicture}
\end{center}
\begin{enumerate}[a)]
    \item Vì $BE, CF$ là các đường cao nên $\widehat{BEC} = \widehat{BFC} = 90^\circ$.\\
    Tứ giác $BCEF$ có hai đỉnh kề $E, F$ cùng nhìn đoạn $BC$ dưới góc $90^\circ$ nên nội tiếp đường tròn đường kính $BC$.\\
    Vì $M$ là trung điểm của $BC$ nên $M$ là tâm đường tròn ngoại tiếp tứ giác $BCEF$, do đó $ME = MF = \dfrac{BC}{2}$.\\
    Xét tứ giác $AEHF$ có $\widehat{AEH} + \widehat{AFH} = 90^\circ + 90^\circ = 180^\circ$ nên nội tiếp đường tròn đường kính $AH$.\\
    Vì $N$ là trung điểm của $AH$ nên $N$ là tâm đường tròn ngoại tiếp tứ giác $AEHF$, do đó $NE = NF = \dfrac{AH}{2}$.\\
    Vì $ME = MF$ và $NE = NF$ nên $MN$ là đường trung trực của đoạn thẳng $EF$. Do đó $MN \perp EF$.
    \item Tứ giác $BCEF$ nội tiếp nên góc ngoài bằng góc đối trong: $\widehat{KFB} = \widehat{KCE}$.\\
    Xét $\triangle KFB$ và $\triangle KCE$ có:
    $$\widehat{EKC}\text{ chung, } \widehat{KFB} = \widehat{KCE}.$$
    Suy ra $\triangle KFB \backsim \triangle KCE$ (g-g) $\implies \dfrac{KF}{KC} = \dfrac{KB}{KE} \implies KB \cdot KC = KF \cdot KE$.
    \item Tứ giác $BCEF$ nội tiếp nên $\widehat{AFE} = \widehat{ACB}$ (cùng bù với $\widehat{BFE}$).\\
    Xét $\triangle AEF$ và $\triangle ABC$ có:
    $$\widehat{BAC}\text{ chung, } \widehat{AFE} = \widehat{ACB}.$$
    Suy ra $\triangle AEF \backsim \triangle ABC$ (g-g) $\implies \dfrac{EF}{BC} = \dfrac{AE}{AB}$.\\
    Trong tam giác vuông $ABE$ tại $E$:
    $$\dfrac{AE}{AB} = \cos \widehat{BAC} = \cos 60^\circ = \dfrac{1}{2}.$$
    Suy ra:
    $$EF = \dfrac{1}{2}BC = \dfrac{1}{2} \cdot 16 = 8\text{ (cm)}.$$
    Vì tứ giác $AEHF$ nội tiếp đường tròn đường kính $AH$ có tâm là $N$ nên đường tròn ngoại tiếp tam giác $AEF$ chính là đường tròn tâm $N$, bán kính $R_{AEF} = NE$.\\
    Gọi $I$ là giao điểm của $MN$ và $EF$. Do $MN$ là đường trung trực của $EF$ nên $I$ là trung điểm của $EF$ và $NI \perp EF$ tại $I$, suy ra:
    $$IE = \dfrac{EF}{2} = \dfrac{8}{2} = 4\text{ (cm)}.$$
    Ta có $\widehat{ENF} = 2\widehat{EAF} = 2 \cdot 60^\circ = 120^\circ$ (góc ở tâm gấp đôi góc nội tiếp cùng chắn cung $EF$).\\
    Tam giác $NEF$ cân tại $N$ có đường cao $NI$, nên $NI$ cũng là phân giác, suy ra $\widehat{INE} = 60^\circ$.\\
    Trong tam giác vuông $NIE$ tại $I$:
    $$\sin \widehat{INE} = \dfrac{IE}{NE} \implies NE = \dfrac{IE}{\sin 60^\circ} = \dfrac{4}{\dfrac{\sqrt{3}}{2}} = \dfrac{8\sqrt{3}}{3}\text{ (cm)}.$$
    Vậy bán kính đường tròn ngoại tiếp tam giác $AEF$ là $\dfrac{8\sqrt{3}}{3}\text{ cm}$.
\end{enumerate}
}
\end{bt}
